\section{Limitações da Revisão}

Cinco limitações delimitam a leitura dos resultados e qualificam o alcance das conclusões.

\noindent\textbf{Granularidade da extração.} A síntese temática deriva dos campos estruturados das fichas (artefatos, papéis, execução, validação, portabilidade, riscos, natureza e evidência) e de um achado narrativo consolidado por estudo, produzido em uma segunda passada de extração. Doze estudos têm achado vindo de uma justificativa de inclusão substantiva e os outros vinte e cinco recebem um achado consolidado a partir dos campos estruturados, sinalizado de forma transparente como tal. Onde a justificativa de primeira passada é genérica, a leitura ainda se apoia em parte na identidade pública conhecida de frameworks bem documentados, como ChatDev, MetaGPT, AutoCodeRover, Reversa e AgileGen. Uma extração de texto completo, campo a campo, para os vinte e cinco estudos restantes aprofundaria ainda mais a síntese.

\noindent\textbf{Avaliação de qualidade por categoria, não por texto completo.} Aplicamos uma avaliação de qualidade de quatro itens por categoria de fonte (Tabela~\ref{tab:qualidade}), que pondera a força da evidência e sinaliza os estudos de menor confiança, mas não substitui um instrumento formal de risco de viés aplicado ao texto completo de cada estudo. Os dez preprints arXiv ainda não revisados por pares e os dois estudos de confiança de extração média entram na síntese com leitura cautelosa, e nenhum tema se apoia exclusivamente neles; ainda assim, uma reanálise de texto completo com um instrumento validado refinaria a ponderação.

\noindent\textbf{Recorte de idioma e de bases, e viés de publicação.} O protocolo restringe a revisão a inglês e português, no período de 2018 a 2026. A busca direta no arXiv falhou em parte por instabilidade de acesso, com a cobertura de preprints recuperada via Semantic Scholar e DOI do arXiv, o que pode ter deixado de fora preprints muito recentes não indexados por essas rotas. Um viés de publicação a favor de frameworks com resultados positivos é provável e mensurável no próprio corpus: 26 dos 37 estudos (70\%) reivindicam avaliação empírica, quase sempre favorável ao framework proposto, enquanto estudos de falha, abandono ou não adoção são praticamente ausentes; apenas o eixo de adoção do tema T5, com seis estudos, problematiza a integração organizacional. Ampliar a busca para bases adicionais e remover a restrição de idioma é trabalho futuro que poderia equilibrar essa assimetria.

\noindent\textbf{Pendência operacional fora do corpus.} Uma fonte aprovada na triagem ficou sem texto completo obtido na data de corte e foi mantida como pendência operacional, fora da extração e desta síntese. O corpus efetivamente sintetizado é de 37 estudos; a incorporação tardia dessa fonte ajustaria contagens marginais, sem alterar a estrutura temática.

\noindent\textbf{Saturação parcial do snowballing.} O snowballing foi encerrado no limite protocolar de três rodadas. A curva de saturação (Figura~\ref{fig:saturacao}) mostra inclusões decrescentes, de 25 na busca inicial para 6, 3 e 3 nas três rodadas: uma desaceleração clara, mas não nula. As últimas rodadas ainda adicionavam estudos, o que indica saturação parcial e não completa; rodadas adicionais poderiam revelar estudos relevantes, em especial fora das comunidades de origem das sementes.
